\section{Method: Mid-Layer Probe Regularization}
\label{sec:method}

MLPR augments parameter-efficient fine-tuning with an observation pathway
(a mid-layer probe and an online recall metric) and a control pathway (a
saturation gate and an eval-aware controller). Figure~\ref{fig:method}
shows the data flow; Algorithm~\ref{alg:mlpr} gives the training loop.

\begin{figure*}[t]
\centering
\resizebox{0.985\textwidth}{!}{%
% MLPR methodology overview — full-width figure* diagram.
% Custom TikZ icons (no external font/icon packages required):
%   snowflake = frozen, flame = trainable, gauge = saturation gate,
%   ramp inset = lambda(t), floor glyph = damping floor, slashed circle = detach.
% Shared by: fig_method.tex (standalone preview) and main paper via \input.
\newcommand{\mini}{\fontsize{5.8}{6.9}\selectfont}
\begin{tikzpicture}[
  x=1cm, y=1cm,
  fwd/.style={-{Stealth[length=5.5pt]}, semithick, black!72},
  tele/.style={-{Stealth[length=5pt]}, semithick, black!55, dashed},
  lamarr/.style={-{Stealth[length=6pt]}, semithick, gateC},
  blk/.style={draw=black!55, rounded corners=1.5pt, fill=gray!10,
              minimum width=0.68cm, minimum height=0.56cm,
              align=center, font=\scriptsize, inner sep=1pt},
  blkhl/.style={blk, fill=frozenBg, draw=frozenC, thick},
  chip/.style={draw=black!55, rounded corners=2pt, align=center,
               inner sep=2.5pt, font=\scriptsize},
  card/.style={chip, fill=white, text width=2.95cm},
  badge/.style={circle, fill=stepC, text=white, font=\tiny\bfseries,
                minimum size=0.34cm, inner sep=0pt},
  mini/.style={font=\mini},
]

% ================= palette =================
\definecolor{frozenC}{RGB}{68,119,170}   % blue  — frozen
\definecolor{frozenBg}{RGB}{232,239,247}
\definecolor{trainC}{RGB}{238,119,51}    % orange — trainable
\definecolor{gateC}{RGB}{170,51,119}     % magenta — gate / lambda
\definecolor{ctrlC}{RGB}{0,153,136}      % teal — controller
\definecolor{gradC}{RGB}{204,51,17}      % red — detach marker
\definecolor{openC}{RGB}{34,136,68}      % green — gate open
\definecolor{evC}{RGB}{238,170,0}        % amber — evidence card
\definecolor{stepC}{RGB}{85,85,85}
\definecolor{anchorBg}{RGB}{253,231,214}

% ================= icons =================
\tikzset{
  flake/.pic={
    \foreach \a in {0,60,120}{\draw[frozenC, line width=0.55pt] (\a-180:0.145) -- (\a:0.145);}
    \foreach \a in {0,60,...,300}{
      \draw[frozenC, line width=0.45pt] (\a:0.085) -- ++(\a+42:0.055);
      \draw[frozenC, line width=0.45pt] (\a:0.085) -- ++(\a-42:0.055);}
    \fill[frozenC] (0,0) circle(0.03);
  },
  flame/.pic={
    \fill[trainC] (0,0.19) .. controls (-0.15,0.04) and (-0.13,-0.07) .. (-0.02,-0.17)
        .. controls (0.10,-0.08) and (0.13,0.02) .. (0,0.19) -- cycle;
    \fill[white] (0,0.055) .. controls (-0.055,-0.02) and (-0.055,-0.075) .. (0,-0.125)
        .. controls (0.055,-0.07) and (0.055,-0.02) .. (0,0.055) -- cycle;
  },
  gauge/.pic={
    \draw[black!55, line width=0.5pt] (-0.30,0) arc[start angle=180, end angle=0, radius=0.30];
    \draw[black!55, line width=0.5pt] (-0.30,0) -- (0.30,0);
    \foreach \ang in {180,135,90,45,0}{\draw[black!45, line width=0.4pt] (\ang:0.26) -- (\ang:0.30);}
    \draw[gateC, line width=1.1pt] (36:0.21) -- (36:0.335);   % tau_0 = 0.8
    \draw[black!70, line width=1.2pt] (0,0) -- (34.9:0.245);  % needle at 0.805
    \fill[black!60] (0,0) circle(0.032);
  },
  ramp/.pic={
    \draw[->, black!55, line width=0.5pt] (-0.05,0) -- (1.42,0);
    \draw[->, black!55, line width=0.5pt] (0,-0.06) -- (0,0.86);
    \draw[gateC, line width=1.1pt] (0,0) -- (0.56,0) -- (1.36,0.76);
    \fill[gateC] (0.56,0) circle(0.035);
    \node[mini, black!60, below] at (0.56,-0.03) {$t_0$};
    \node[mini, black!60, left]  at (-0.03,0.74) {$\lambda_0$};
  },
  floorplt/.pic={
    \draw[->, black!55, line width=0.5pt] (-0.05,0.72) -- (-0.05,0);
    \draw[->, black!55, line width=0.5pt] (-0.05,0) -- (1.52,0);
    \draw[ctrlC, line width=1.1pt] (0,0.62) -- (0.30,0.40) -- (0.58,0.27) -- (0.82,0.21)
        -- (1.05,0.19) -- (1.45,0.19);
    \draw[black!45, dashed, line width=0.5pt] (-0.05,0.19) -- (1.45,0.19);
    \node[mini, black!60, below] at (1.18,0.0) {$m_{\min}$};
    \node[mini, black!60, left]  at (-0.05,0.62) {$m$};
  },
  det/.pic={ % detach: slashed circle
    \fill[white] (0,0) circle(0.085);
    \draw[gradC, line width=0.8pt] (0,0) circle(0.085);
    \draw[gradC, line width=0.8pt] (135:0.115) -- (-45:0.115);
  },
}

% ================= band backgrounds =================
\fill[blue!3, rounded corners=4pt]   (0.10,-1.50) rectangle (16.98,3.48);
\fill[black!3, rounded corners=4pt]  (0.10,-6.02) rectangle (16.98,-1.72);
\node[anchor=west, font=\scriptsize\bfseries, black!55] at (0.32,3.26)
  {DATA \& OBSERVATION PATHWAY \; (per batch)};
\node[anchor=west, font=\scriptsize\bfseries, black!55] at (0.32,-1.78)
  {CONTROL PATHWAY \; (per epoch)};

% ================= worked-example query card =================
\node[card, anchor=north west] at (0.35,2.62) {%
  {\scriptsize\bfseries Worked example}\\[2pt]
  {\mini\itshape Q:} Which regulatory authority oversees
  \colorbox{anchorBg}{\mini Institution\_Alpha\_175}?\\[3pt]
  {\mini\itshape A:} \texttt{Answer:}\\[0.5pt]
  {\mini\hspace*{1.05em} Regulator\_Beta\_044}\\[3pt]
  {\mini probe label: entity id 243}\\
  {\mini anchor token: chars 36--57}};

% ================= backbone chain =================
\node[circle, draw=black!55, fill=gray!8, minimum size=0.42cm,
      font=\scriptsize, inner sep=0pt] (emb) at (4.12,1.9) {$x$};
\node[blk] (L1)  at (4.94,1.9) {$1$};
\node[blk] (L2)  at (5.78,1.9) {$2$};
\node[font=\scriptsize, black!55] at (6.42,1.9) {$\cdots$};
\node[blkhl] (LS) at (7.06,1.9) {$l^\star$};
\node[blk] (LL)  at (8.34,1.9) {$L$};
\pic at (4.98,2.80) {flake};
\node[anchor=west, font=\mini, black!60] at (5.18,2.80)
  {decoder layers $1\ldots L$ --- backbone $\theta$ \textbf{frozen}};

% LoRA branch (trainable)
\node[chip, fill=trainC!10, draw=trainC!80!black, minimum width=1.02cm] (Abox) at (5.00,-0.90)
  {\scriptsize LoRA $A$\\{\mini $d\!\to\!r{=}16$}};
\node[chip, fill=trainC!10, draw=trainC!80!black, minimum width=1.02cm] (Bbox) at (6.42,-0.90)
  {\scriptsize LoRA $B$\\{\mini $r\!\to\!d$}};
\pic at (4.66,-0.615) {flame};
\pic at (6.08,-0.615) {flame};
\node[font=\mini, black!60, align=center, text width=3.0cm] at (5.62,-0.22)
  {$\Delta W{=}\alpha r^{-1}BA$ per layer\\(7 projections: attn + MLP)};
\draw[fwd] (emb.south) -- (4.12,-0.90) -- (Abox.west);
\draw[fwd] (Abox.east) -- (Bbox.west);
\node[circle, draw=black!60, fill=white, inner sep=0.4pt, font=\tiny] (plus) at (7.62,1.9) {$+$};
\draw[fwd] (Bbox.east) -- (7.62,-0.90) -- (plus.south);

% LM head / answer / L_ce
\node[chip, fill=gray!8, minimum width=0.92cm] (head) at (9.32,1.9) {\scriptsize LM head};
\node[chip, fill=white, minimum width=2.42cm] (ans) at (11.55,1.9)
  {\scriptsize Answer: Regulator\_Beta\_044};
\node[chip, fill=blue!8, minimum width=1.12cm, minimum height=0.72cm] (lce) at (14.00,1.9)
  {\scriptsize $\mathcal{L}_{\mathrm{ce}}$\\{\mini SFT loss}};
\draw[fwd] (emb.east) -- (L1.west);
\draw[fwd] (LL.east) -- (head.west);
\draw[fwd] (head.east) -- (ans.west);
\draw[fwd] (ans.east) -- (lce.west);

% epoch-end eval node
\node[chip, fill=gray!6, text width=2.0cm] (evaln) at (15.75,2.94)
  {\scriptsize\bfseries epoch-end eval\\{\mini 256-prompt recall subsample}\\
   {\mini $A_{\mathrm{mem}}^{\mathrm{lf}},\,A_{\mathrm{mem}}^{\mathrm{em}},\,A_{\mathrm{gen}}$}};
\draw[tele] (ans.north) |- (evaln.west);

% probe pathway
\node[chip, fill=trainC!10, draw=trainC!80!black, text width=2.52cm] (probe) at (6.05,0.50)
  {\scriptsize\bfseries linear probe $p_\phi$\\{\mini $d\!\to\!|V_e|{=}1600$ · LR $2{\times}10^{-3}$}};
\pic at (4.62,0.82) {flame};
\node[chip, fill=white, minimum width=2.26cm] (pred) at (9.00,0.50)
  {\scriptsize entity: Regulator\_Beta\_044\\{\mini probe argmax = id 243}};
\node[chip, fill=trainC!10, draw=trainC!80!black, minimum width=1.30cm,
      minimum height=0.72cm] (lpr) at (11.42,0.50)
  {\scriptsize $\mathcal{L}_{\mathrm{probe}}$\\{\mini entity CE}};
\draw[fwd] (LS.south) -- (7.06,0.84);
\node[mini, black!60, anchor=east] at (6.94,1.12) {$h_{l^\star}$ (anchor)};
\pic at (7.06,1.42) {det};
\node[mini, gradC, anchor=east] at (6.90,1.42) {detached (warm-up)};
\draw[fwd] (probe.east) -- (pred.west);
\draw[fwd] (pred.east) -- (lpr.west);

% ================= control pathway =================
% gate
\node[chip, fill=gateC!8, draw=gateC, thick, anchor=west, text width=2.42cm,
      minimum height=1.24cm] (gate) at (0.30,-2.70)
  {\scriptsize\bfseries saturation gate\\[2pt]
   {\mini $\rho{\in}\{A_{\mathrm{mem}}^{\mathrm{lf}},\,\rho_{\mathrm{rep}}\}$}\\[1pt]
   {\mini $\ge\ \tau_0$\;?}};
\pic at (3.30,-2.62) {gauge};

% closed / open chips
\node[chip, fill=gray!8, minimum width=1.52cm] (closed) at (0.95,-3.74)
  {\mini closed: $\lambda{=}0$\\{\mini pure SFT}};
\node[chip, fill=openC!10, draw=openC!80!black, minimum width=1.56cm] (open) at (2.72,-3.74)
  {\mini open at $t_0$:\\{\mini ramp $\lambda(t)$}};
\draw[fwd] (gate.south) -- ++(0,-0.16) -| (closed.north);
\draw[fwd] (gate.south) -- ++(0,-0.16) -| (open.north);

% evidence card
\node[chip, fill=evC!12, draw=evC!75!black, text width=3.42cm] (ev) at (1.95,-4.94)
  {\scriptsize\bfseries worked example --- 1.7B run\\[1.5pt]
   {\mini ep 17: $A_{\mathrm{mem}}^{\mathrm{lf}}{=}0.885 \ge \tau_0$}\\[0.5pt]
   {\mini gate opens at $\lambda_0{=}0.3$}\\[0.5pt]
   {\mini ep 18: signal re-enters band}\\[0.5pt]
   {\mini $\lambda$: $0.3{\to}0.075{\to}0.3$ (live)}};

% lambda ramp inset
\pic at (4.30,-2.62) {ramp};
\node[font=\mini, black!60, align=center, anchor=north, text width=2.6cm] at (4.45,-2.92)
  {$\lambda(t){=}\lambda_0\,\min\!\bigl(1,\tfrac{\rho(t)-\tau_0}{\Delta_0}\bigr)$\\[1pt]
   {\mini $\lambda_0{=}0.3$; decay variant:}\\[0.5pt]
   {\mini $\times\max(0,1{-}\tfrac{t-t_0}{D})$}};

% total loss
\node[chip, fill=yellow!14, draw=black!60, thick, minimum width=2.56cm,
      minimum height=1.18cm] (loss) at (8.15,-2.62)
  {\scriptsize\bfseries total objective\\[1.5pt]
   {\scriptsize $\mathcal{L}=\mathcal{L}_{\mathrm{ce}}+\lambda(t)\,\mathcal{L}_{\mathrm{probe}}$}};
\draw[lamarr] (5.80,-2.62) -- (6.87,-2.62);
\node[badge] at (6.33,-2.62) {5};
\node[mini, gateC, anchor=north] at (6.33,-2.84) {$\lambda(t)$};

% controller + optimizer
\node[chip, fill=ctrlC!8, draw=ctrlC, thick, text width=2.92cm] (ctrl) at (12.50,-2.62)
  {\scriptsize\bfseries eval-aware controller\\[2pt]
   {\mini boost LR while held-out improves;}\\[0.5pt]
   {\mini decay $\times\,0.7$ after degradation;}\\[0.5pt]
   {\mini floor $m\leftarrow\max(0.5,\,m\gamma)$}};
\pic at (12.50,-4.18) {floorplt};
\node[circle, draw=black!60, fill=gray!10, minimum size=1.06cm, font=\tiny, inner sep=0pt]
  (opt) at (15.42,-2.62) {\textbf{AdamW}};
\draw[fwd] (loss.east) -- node[above, mini, black!60] {gradients} (ctrl.west);
\draw[fwd] (ctrl.east) -- node[above, mini, black!60] {LR $\times\,m(t)$} (opt.west);

% trainable scope chip
\node[chip, fill=trainC!8, draw=trainC!80!black, text width=1.9cm] (scope) at (15.42,-3.95)
  {\mini \tikz{\pic[scale=0.62]{flame}}\;updates only $A,B,\phi$\\[1pt]
   \tikz{\pic[scale=0.62]{flake}}\;$\theta$ stays frozen};
\draw[fwd] (opt.south) -- (scope.north);

% ================= telemetry / loss feeds =================
\draw[tele] (10.60,1.64) -- (10.60,-1.55) -| (1.60,-1.97);
\node[badge] at (10.60,-0.36) {4};
\draw[fwd] (lce.south) -- (14.00,-1.32) -- (8.72,-1.32) -- (8.72,-2.03);
\draw[fwd] (lpr.south) -- (11.42,-1.14) -- (7.62,-1.14) -- (7.62,-2.03);

% ================= legend =================
\node[chip, fill=white, text width=2.82cm, anchor=north] (leg) at (15.42,-4.50) {%
  {\mini\bfseries legend}\\[1.5pt]
  {\mini \tikz{\pic[scale=0.62]{flake}} frozen $\theta$ \;\; \tikz{\pic[scale=0.62]{flame}} trainable}\\[1pt]
  {\mini \tikz{\draw[black!55,semithick,dashed,-{Stealth[length=3pt]}](0,0)--(0.42,0);} eval-only signal}\\[1pt]
  {\mini \tikz{\pic[scale=0.62]{det}} detached (no grad)}};

% ================= walkthrough badges =================
\node[badge] at (3.70,1.9) {1};
\node[badge] at (7.06,2.44) {2};
\node[badge] at (7.06,1.10) {3};
\node[badge] at (15.42,-3.36) {6};

\end{tikzpicture}
}
\caption{MLPR methodology with a worked example.
\textbf{Top --- data \& observation pathway (per batch):} a memorization
sample from $\mathcal{D}_{\mathrm{mem}}$ (left card; anchor token
highlighted) is processed by the frozen backbone (snowflake: $\theta$
receives no gradient) augmented with trainable low-rank adapters
(flame: $A,B$, $r{=}16$, applied to all seven attention and MLP
projections, merged at $\oplus$). The mid-layer state $h_{l^\star}$ at
the anchor token is read by a linear probe $p_\phi$ (flame) that
predicts the entity identity; during warm-up the probe observes a
\emph{detached} copy of the state ($\oslash$, no gradient into
$\theta$). \textbf{Bottom --- control pathway (per epoch):} the
saturation gate (gauge) compares teacher-forced recall
$\rho=\max(\amemlf,\amemem)$ against $\tau_0{=}0.8$; once it opens,
$\lambda(t)$ ramps the probe-alignment loss into the total objective
(ramp inset), and an eval-aware controller modulates the learning rate
under a damping floor $m\ge m_{\min}{=}0.5$ (floor inset). Numbered
badges trace the example: (1)~sample enters; (2)~frozen+LoRA forward;
(3)~detached mid-layer probe read; (4)~epoch-end recall evaluation
(208-prompt subsample); (5)~gated loss weight; (6)~optimizer updates
trainable parameters only. The amber card reports the realized 7B run:
the gate opened at epoch 24 ($\amemlf{=}0.805\ge\tau_0$) and $\lambda$
ramped $0.015\!\to\!0.075\!\to\!0.225$ toward $\lambda_0{=}0.3$,
whereas the 0.5B run never crossed the threshold ($\max 0.39$).}
\label{fig:method}
\end{figure*}

\subsection{Setup and objective}
\label{sec:method:setup}

Let $f_\theta$ be a pretrained decoder-only LM, and let LoRA
\citep{hu2022lora} add low-rank updates
$\Delta W = \tfrac{\alpha}{r} BA$ to selected weight matrices, with $A,B$
trainable and $\theta$ frozen. The supervised objective on the
memorization corpus $\mathcal{D}_{\mathrm{mem}}$ of question--answer
pairs is the token-level cross-entropy
\begin{equation}
  \mathcal{L}_{\mathrm{ce}}
  = -\frac{1}{|B|}\sum_{(x,y)\in B}
    \log p_\theta\!\bigl(y \mid x\bigr),
  \label{eq:ce}
\end{equation}
where $x$ is the training-format question and $y$ the answer with
\texttt{Answer:} suffix. At each optimizer step the total loss is
\begin{equation}
  \mathcal{L} = \mathcal{L}_{\mathrm{ce}}
  + \lambda(t)\, \mathcal{L}_{\mathrm{probe}}\bigl(p_\phi, f_\theta\bigr),
  \label{eq:total}
\end{equation}
where $\lambda(t)\!\in\![0,\lambda_0]$ is the gate weight
(Section~\ref{sec:method:gate}) and $\mathcal{L}_{\mathrm{probe}}$ is a
cross-entropy between the probe's entity prediction and the ground-truth
entity id of the sample in the batch, computed on the mid-layer hidden
state $h_{l^\star}$ \emph{with gradients flowing into the backbone} once
$\lambda>0$ (before activation the probe receives no gradient path into
$\theta$, cf.\ Eq.~\ref{eq:total} with $\lambda{=}0$).

\subsection{Probe warm-up and the likelihood recall metric}
\label{sec:method:probe}

The probe $p_\phi:\mathbb{R}^{d}\!\to\!\mathbb{R}^{V_e}$ is a single
linear layer over the entity vocabulary $V_e$ (here $|V_e|{=}1600$),
trained from the first step on detached states at
$l^\star{\approx}L/2$ with its own learning rate
($10^{-3}$ vs.\ $2{\times}10^{-5}$ for the backbone), so that it
converges quickly enough to be informative early. Because the probe
gradient is detached during warm-up, it is a pure observer; only after
gate activation does $\mathcal{L}_{\mathrm{probe}}$ back-propagate into
$\theta$, at which point it encourages the backbone to keep acquired
entity identities linearly decodable at $l^\star$.

Free-generation exact match is a poor saturation signal: a model may
encode the correct answer yet fail to emit it in the right surface form.
We therefore measure recall with a teacher-forced likelihood probe: for
each memorization prompt (answer suffix stripped), the answer tokens are
scored in a single forward pass and the greedy per-sample match is
\begin{equation}
  \mathrm{lf}(x, y) = \mathbb{1}\bigl[
    \arg\max_w p_\theta(w \mid x) = y_1 \bigr],
  \label{eq:lf}
\end{equation}
with a token-level variant that credits partial answers. The epoch-level
recall is $\amemlf = \mathbb{E}_{\mathcal{S}}[\mathrm{lf}]$ on a fixed
subsample $\mathcal{S}$ of $\mathcal{D}_{\mathrm{mem}}$ (208 prompts).
We additionally log the free-generation exact match
$\amemem$, substring containment, paraphrased generalization accuracy
$\agen$ on $\mathcal{D}_{\mathrm{gen}}$, and the probe's own
train/held-out accuracies, giving a five-channel view of the same
underlying state change.

\subsection{Saturation gate and $\lambda(t)$}
\label{sec:method:gate}

Let $t$ index the training step and $T$ the total step budget. The gate
uses the maximum of the likelihood and free-generation recall signals,
$\rho(t) = \max(\amemlf, \amemem)$, and a threshold $\tau_0$. Before
$\rho \ge \tau_0$ the gate is closed ($\lambda{=}0$) and training is
pure SFT. After the first crossing at step $t_0$, the weight ramps
linearly over a fraction $\delta_0$ of the budget,
\begin{equation}
  \lambda(t) = \lambda_0 \cdot
    \min\!\Bigl(1,\ \tfrac{t - t_0}{\delta_0 T}\Bigr),
  \qquad t \ge t_0,
  \label{eq:ramp}
\end{equation}
so the regularization is introduced gradually ($\lambda_0{=}0.3$,
$\tau_0{=}0.8$, $\delta_0{=}0.1$ here). The maximum over the two recall
modes makes the gate robust to the answer-extraction artifact
\citep{mallen2023pararel}: in pilot runs the free-generation signal
plateaued at $0$--$12\%$ and never opened a gate that required it, while
the likelihood signal crossed the same threshold at scale
(Section~\ref{sec:res:sat}).

\subsection{Eval-aware adaptive control with a damping floor}
\label{sec:method:control}

Long memorization runs exhibit a characteristic trajectory: training
loss plateaus after the format is learned, and held-out loss---bottoming
early---then rises monotonically as capacity is redirected to training
entities \citep{tirumala2022memorization}. An adaptive controller that
treats the plateau as under-optimization may boost the learning rate
exactly while validation degrades, and one that treats the rise as
divergence may decay it without bound. Both failure modes occurred in
our pilot runs. The controller we ship therefore (i) boosts only while
held-out loss is not degrading, (ii) decays by a fixed factor after
repeated degradation episodes, and (iii) enforces a \emph{damping floor}:
the cumulative LR multiplier $m(t)$ applied by eval-driven decays is
clamped to $m(t) \ge m_{\min}$ with $m_{\min}{=}0.5$,
\begin{equation}
  m(t) \leftarrow \max\!\bigl(m_{\min},\ m(t^-)\cdot \gamma\bigr),
  \label{eq:floor}
\end{equation}
with decay factor $\gamma{=}0.7$. The floor is the load-bearing fix: in
the memorization regime, held-out divergence is the phenomenon under
study, not a fault, and unbounded damping freezes it
(Section~\ref{sec:res:control} quantifies the spiral we observed without
the floor). Every decision is logged per epoch, and all controller
actions are failure-isolated so that instrumentation cannot crash
training.

\begin{algorithm}[t]
\caption{MLPR training loop (one epoch)}
\label{alg:mlpr}
\KwIn{corpora; LoRA backbone $f_\theta$; probe $p_\phi$; gate
  \eqref{eq:ramp}; controller with floor \eqref{eq:floor}}
probe warm-up on detached $h_{l^\star}$\;
\ForEach{batch $B \subset \mathcal{D}_{\mathrm{mem}}$}{
  $\mathcal{L} \leftarrow \mathcal{L}_{\mathrm{ce}}
    + \lambda(t)\,\mathcal{L}_{\mathrm{probe}}$; update $\{A,B,\phi\}$\;
}
$\amemlf \leftarrow$ \eqref{eq:lf} on subsample; evaluate
$\amemem, \agen$\;
open gate if $\rho \ge \tau_0$; clamp controller $m \ge m_{\min}$; log\;
\end{algorithm}
